\section{Appendix}
\label{sec:appendix}


\subsection{OpenBrew Platform Data Collection}
\label{sec:openbrew_platform_data_collection}

\begin{lstlisting}[caption={Example of collected platform data}, label={lst:openbrew_data}, basicstyle=\ttfamily\footnotesize, frame=single, backgroundcolor=\color{gray!10}, breaklines=true]
{
    "users": [
        {
            "id": "user123",
            "username": "coffee_lover",
            "role": "enthusiast"
        },
        {...},
        {
            "id": "user101",
            "username": "roastmaster",
            "role": "roaster"
        }
    ],
    "posts": [
        {
            "id": "post456",
            "user_id": "user123",
            "content": "Just brewed a delicious cup of coffee!",
            "timestamp": "2024-04-01T12:00:00Z",
            "upvotes": 10,
            "downvotes": 2,
            "upvoted_by": ["user456", ..., "user789"],
            "downvoted_by": ["user101"],
            "comments": [
                {
                    "user_id": "user456",
                    "content": "What beans did you use?",
                    "timestamp": "2024-04-01T12:05:00Z",
                    "upvotes": 3,
                    "downvotes": 0,
                    "upvoted_by": [...],
                    "downvoted_by": [],
                    "comments": [
                        {
                            "user_id": "user789",
                            "content": "I bet it was Ethiopian!",
                            "timestamp": "2024-04-01T12:06:00Z",
                            "upvotes": 1,
                            ...,
                            "comments": []
                        }
                    ]
                }
            ]
        }
    ]
}
\end{lstlisting}

\newpage
\begin{lstlisting}[caption={Example output from the basic statistics API}, label={lst:openbrew_basic_stats}, basicstyle=\ttfamily\footnotesize, frame=single, backgroundcolor=\color{gray!10}, breaklines=true]
{
  "start_time": "2024-01-01T00:00:00Z",
  "end_time": "2024-01-31T23:59:59Z",
  "number_of_users": 152,
  "number_of_brews": 8,
  "number_of_posts": 412,
  "number_of_comments": 1297
}
\end{lstlisting}

\subsection{Pseudocode for Constructing Network Graphs}
\label{sec:pseudocode_for_constructing_network_graphs}

\subsubsection{Votes Graph Construction}
\begin{lstlisting}[language=Python, caption={Pseudocode for Votes Graph Construction}, basicstyle=\ttfamily\footnotesize, frame=single, backgroundcolor=\color{gray!10}, breaklines=true]
For each user in users:
    Add node to graph with user id and role

For each post in posts:
    For each user in post.upvoted_by:
        If user != post.user_id:
            Add or update directed edge (user -> post.user_id) with +1 weight and timestamp
    For each user in post.downvoted_by:
        If user != post.user_id:
            Add or update directed edge (user -> post.user_id) with -1 weight and timestamp
    Recursively process comments as posts
\end{lstlisting}

\subsubsection{Comments Graph Construction}
\begin{lstlisting}[language=Python, caption={Pseudocode for Comments Graph Construction}, basicstyle=\ttfamily\footnotesize, frame=single, backgroundcolor=\color{gray!10}, breaklines=true]
For each user in users:
    Add node to graph with user id and role

For each post in posts:
    Recursively process comments:
        For each comment:
            Add or update directed edge (comment.user_id -> parent_author_id) with +1 weight and timestamp
            Recursively process nested comments (parent_author_id = comment.user_id)
\end{lstlisting}
